\section{Relative polynomial generators and the original envelope}

\begin{definition}[Bounded polynomial cone generators]
\label{def:c0n-bounded-polynomial-cone-generators}
\uses{conv:sg-setting,lem:sg-phi-basic}
\lean{EnvelopeProperty.ProjSubvariety.HasBoundedPolynomialConeGenerators}
\leanok
An ideal $B\subset T$ has bounded polynomial cone generators of degree $e$
if there is a finite set $G\subset k[x_0,\ldots,x_M][t]$ with
\[
 (\bar G)\subset B\subset(\bar G)^{\mathrm{sat}}
\]
such that every coefficient of every $F\in G$, viewed as a polynomial in
$t$, is a homogeneous form in the $x$ variables of some degree at most $e$.
Here $\bar G$ is the image in the actual cone ring $T$.
\end{definition}

\begin{lemma}[Fubini--Study functions from relative polynomial generators]
\label{lem:c0n-polynomial-cone-fs}
\uses{def:c0n-bounded-polynomial-cone-generators,lem:sg-phi-basic}
\lean{EnvelopeProperty.ProjSubvariety.isFS_coneIdealApprox_of_polynomial_generators}
\leanok
If $B_m\ne0$ has bounded polynomial cone generators of degree $md+c$,
where $m>0$, then $m^{-1}\phi_{B_m}\in\mathcal H(X,d+c/m)$.
\end{lemma}
\begin{proof}
Let $G$ be the finite polynomial set in the definition.  Saturation does
not change its potential on $X^{\mathrm{an}}$, and adding the kernel of
$k[x_0,\ldots,x_M,t]\to T$ does not change it there either.  Thus
$\phi_{B_m}=\phi_{(G)}$ on $X^{\mathrm{an}}$.
Some $F\in G$ has nonzero image in $T$: otherwise $(\bar G)=0$, whose
saturation is zero since $T$ is a domain and its irrelevant ideal contains
a nonzero coordinate, contradicting $B_m\ne0$.
The coefficient degree bound and the finite-generator formula for ideal
potentials give $\phi_{(G)}\in\mathcal H(X,md+c)$.
Multiplying this Fubini--Study function by $m^{-1}$ changes the class to
$d+c/m$, as asserted.
\end{proof}

\begin{theorem}[Envelope assembly from relative polynomial generators]
\label{thm:c0n-polynomial-cone-envelope-assembly}
\source{bj22-trivval:thm:contenvlisse}
\uses{lem:c0n-polynomial-cone-fs,lem:c0n-cone-ideal-bounds,thm:pg-psh2,thm:sg-comparison}
\lean{EnvelopeProperty.ProjSubvariety.isPSH_envelope_of_polynomialConeIdealEstimates}
\leanok
Assume that $k$ is algebraically closed, let $d>0$, and let $\phi\le0$ be
piecewise linear on the original $X^{\mathrm{an}}$.
Suppose that actual cone ideals $B_m$ contain $J_m(\phi)T$ for $m>0$,
that $B_{ml}\subset(B_m^l)^{\mathrm{sat}}$ for $m,l>0$, and that for one
fixed $c\in\mathbf N$ every $B_m$ with $m>0$ has bounded polynomial
cone generators of degree $md+c$.
Suppose also that at every divisorial point $v$ there is $A_v\ge0$ such
that for every $m>0$ some $q>0$ satisfies
\[
 \phi_{B_m}(v)\le(m/q)\phi_{J_q(\phi)}(v)+A_v.
\]
Then the original envelope $P_d(\phi)$ is $d$-plurisubharmonic.
\end{theorem}
\begin{proof}
Bound $\phi$ below by a real constant $C$, using continuity and compactness.
The inclusion of the nonzero obstacle ideal gives $B_m\ne0$ for $m>0$.
Put $u_m=m^{-1}\phi_{B_m}$ and $U=\inf_j u_{j!}$.
The preceding lemma gives $u_m\in\operatorname{PSH}(X,d+c/m)$.
The ideal bounds give $u_{(j+1)!}\le u_{j!}$ and
$U\ge P_d(\phi)\ge C$; the second inequality holds because $C$ is a
competitor.  The classes $d+c/j!$ converge to $d$, so the decreasing-limit
theorem gives $U\in\operatorname{PSH}(X,d)$.
At a divisorial point the discrepancy estimate gives
$U(v)\le P_d(\phi)(v)+A_v/j!$ for every $j$.
Given $\varepsilon>0$, choose $j>A_v/\varepsilon$ and use $j!\ge j$.
It follows that $U(v)\le P_d(\phi)(v)\le\phi(v)$.
Divisorial comparison with the continuous obstacle gives $U\le\phi$
everywhere.  Thus $U$ is a competitor for the envelope and
$U\le P_d(\phi)$, whereas the reverse inequality was already established.
Consequently $P_d(\phi)=U$ is $d$-plurisubharmonic.
\end{proof}

\begin{remark}
Relative sections over $k[t]$ yield homogeneous polynomials in the $x$
variables with polynomial coefficients in $t$.  The criterion retains
these polynomials as generators.  It requires no assertion that their
individual $t$-homogeneous components belong to the multiplier ideal.
\end{remark}
